% Part: incompleteness
% Chapter: incompleteness-provability
% Section: second-incompleteness-thm

\documentclass[../../../include/open-logic-section]{subfiles}

\begin{document}

\olfileid{inc}{inp}{2in}

\olsection{第二不完全性定理}

我们该如何表达“$\Th{PA}$ 不能证明自身一致性”这一断言呢？说 $\Th{PA}$ 不一致，就等于说 $\Th{PA} \Proves \eq[0][1]$。因此，我们可以将一致性语句 $\OCon[\Th{PA}]$ 取为语句 $\lnot
\OProv[\Th{PA}](\gn{\eq[0][1]})$，那么下面的定理便说明了这一点：

\begin{thm}
\ollabel{thm:second-incompleteness}
若 $\Th{PA}$ 是一致的，则 $\Th{PA}$ 不能推导 $\OCon[\Th{PA}]$。
\end{thm}

需要注意的是，该定理依赖于 $\OCon[\Th{PA}]$ 的具体表示（即 $\OProv[\Th{PA}](y)$ 的具体表示）。我们只需用到 $\OProv[\Th{PA}](y)$ 的表示满足三个可证性条件这一事实，因此该定理可推广至任何具有满足这些性质的可证性谓词的理论。

阅读 Gödel 的论证概要很有教益，因为定理的得出就像绝妙的点睛之笔一样。论证过程如下。设 $!G_{\Th{PA}}$ 是我们在 \olref[1in]{thm:first-incompleteness} 的证明中构造的 Gödel 语句。我们已经证明了“如果 $\Th{PA}$ 是一致的，那么 $\Th{PA}$ 不能推导 $!G_{\Th{PA}}$”。如果我们将此结论\emph{在} $\Th{PA}$ 中形式化，就得到了 \[
\OCon[\Th{PA}] \lif \lnot \OProv[\Th{PA}](\gn{!G_{\Th{PA}}}).
\] 的证明。现在假设 $\Th{PA}$ 推导 $\OCon[\Th{PA}]$。那么它就推导 $\lnot
\Prov[\Th{PA}](\gn{!G_{\Th{PA}}})$。但由于 $!G_{\Th{PA}}$ 是一个 Gödel 语句，这就等价于 $!G_{\Th{PA}}$。所以 $\Th{PA}$ 推导 $!G_{\Th{PA}}$。

但是：我们知道如果 $\Th{PA}$ 是一致的，它就不能推导 $!G_{\Th{PA}}$!{} 因此，如果 $\Th{PA}$ 是一致的，它就不能推导 $\OCon[\Th{PA}]$。

为了使论证更精确，我们将令 $!G_{\Th{PA}}$ 是 $\Th{PA}$ 的 Gödel 语句，并利用可证性条件 (P1)--(P3) 来证明 $\Th{PA}$ 推导 $\OCon[\Th{PA}] \lif
!G_{\Th{PA}}$。这将表明 $\Th{PA}$ 不能推导 $\OCon[\Th{PA}]$。以下是在~$\Th{PA}$ 中进行的证明概要。（为简单起见，我们省略 $\Th{PA}$ 的下标。）
\begin{align}
& !G \liff \lnot \OProv(\gn{!G}) \ollabel{G2-1}\\
& \qquad\text{$!G$ 是一个 Gödel 语句}\notag \\
& !G \lif \lnot \OProv(\gn{!G}) \ollabel{G2-2}\\
  & \qquad\text{由 \olref{G2-1}} \notag\\
& !G \lif
  (\OProv(\gn{!G}) \lif \lfalse) \ollabel{G2-3}\\
  & \qquad\text{由 \olref{G2-2}，据逻辑}\notag\\
& \OProv(\gn{
    !G \lif
    (\OProv(\gn{!G}) \lif \lfalse)
  }) \ollabel{G2-4}\\
  & \qquad\text{由 \olref{G2-3}，据条件 P1} \notag\\
& \OProv(\gn{!G}) \lif
  \OProv(\gn{
    (\OProv(\gn{!G}) \lif \lfalse)
    }) \ollabel{G2-5}\\
  & \qquad\text{由 \olref{G2-4}，据条件 P2} \notag\\
& \OProv(\gn{!G}) \lif (\OProv(\gn{\OProv(\gn{!G})}) \lif \OProv(\gn{\lfalse})) \ollabel{G2-6}\\
  & \qquad\text{由 \olref{G2-5}，据条件 P2 与逻辑} \notag\\
& \OProv(\gn{!G}) \lif
  \OProv(\gn{\OProv(\gn{!G})}) \ollabel{G2-7}\\
   & \qquad\text{据 P3} \notag\\
& \OProv(\gn{!G}) \lif \OProv(\gn{\lfalse}) \ollabel{G2-8}\\
  & \qquad \text{由 \olref{G2-6} 和 \olref{G2-7}，据逻辑}\notag\\
& \OCon \lif \lnot \OProv(\gn{!G}) \ollabel{G2-9}\\
  & \qquad\text{\olref{G2-8} 的逆否命题及 $\OCon \ident \lnot \OProv(\gn{\lfalse})$}\notag \\
& \OCon \lif !G \notag\\
  & \qquad\text{由 \olref{G2-1} 与 \olref{G2-9}，据逻辑}\notag
\end{align}
上述逻辑的运用仅涉及命题逻辑的基本事实，例如，\olref{G2-3} 用到了 $\Proves \lnot!A \liff (!A\lif
\lfalse)$，而 \olref{G2-8} 用到了 $!A \lif (!B \lif !C), !A \lif !B
\Proves !A \lif !C$。在 \olref{G2-5} 和 \olref{G2-6} 中对条件~P2 的运用依赖于~P2 的实例，$\OProv(\gn{!A \lif !B}) \lif
(\OProv(\gn{!A}) \lif \OProv(\gn{!B}))$。在第一个中，$!A \ident !G$ 且 $!B \ident \OProv(\gn{!G}) \lif \lfalse$；在第二个中，$!A \ident \OProv(\gn{G})$ 且 $!B \ident \lfalse$。

第二不完全性定理更抽象的版本如下：

\begin{thm}
\ollabel{thm:second-incompleteness-gen} 设 $\Th{T}$ 是任意一致的、扩充 $\Th{Q}$ 的可公理化理论，且 $\OProv[\Th{T}](y)$ 是满足 $\Th{T}$ 的可证性条件 P1--P3 的任意公式。则 $\Th{T}$ 不能推导~$\OCon[T]$。
\end{thm}

\begin{prob}
证明 $\Th{PA}$ 推导 $!G_{\Th{PA}} \lif \OCon[\Th{PA}]$。
\end{prob}

\begin{digress}
这里的启示是，没有任何“合理的”一致数学理论能够推导出其自身的一致性语句。假设 $\Th{T}$ 是一个包含 $\Th{Q}$ 和 Hilbert “有穷”推理（无论那可能是什么）的数学理论。那么，整个 $\Th{T}$ 无法推导 $\Th{T}$ 的一致性语句，因此，其有穷片段更是无法推导~$\Th{T}$ 的一致性语句。在这个意义上，不可能存在针对“全部数学”的有穷一致性证明。

在解释“有穷”这个术语时有一定的灵活性，Gödel 在 1931 年的论文中也承认，可能存在某些我们认为是“有穷”的东西，但它们却超出了 Hilbert 想要形式化的数学范围。但 Gödel 的态度是宽容的；如今，很难想象我们会找到某种可以合理地称为“有穷”、却又无法在（比如说）带选择公理的 Zermelo--Fraenkel 集合论 $\Th{ZFC}$ 中被形式化的东西。
\end{digress}

\end{document}
